\begingroup\fontsize{10}{11.8}\selectfont\setlength{\parskip}{2pt}
\section{Conclusions}
\label{v:conclusiones}
Les distributions de Bose--Einstein et de Fermi--Dirac s’obtiennent comme lois
d’équilibre des complétions symétrique et extérieure des états
avec trajectoire et mémoire. L’exclusion algébrique de Pauli et la limite
diluée de Maxwell--Boltzmann font partie de cette même dérivation.
La réalisation électromagnétique de l’occupation bosonique produit,
ensuite, le spectre de Planck, le flux de Stefan--Boltzmann et
les lois de déplacement de Wien. La connexion est ainsi établie
entre la structure de l’état, les occupations admissibles et les observables
radiatifs dans les représentations spécifiées.

Le mécanisme statistique est déterminé par la parité compositionnelle.
Le produit extérieur rend nilpotents les opérateurs de création et
idempotents les opérateurs d’occupation, avec un spectre contenu dans
\(\{0,1\}\) ; la complétion symétrique admet des occupations entières non
négatives. La seconde quantification transporte les entrelaceurs
contractifs, avec le domaine supplémentaire indiqué pour les morphismes
hors de cette classe. L’identification relativiste de la parité et du spin
maintient les conditions de représentation des rotations et de localité :
l’exclusion algébrique et cette identification ont des hypothèses distinctes.

La factorisation grand canonique transforme ces occupations en les
dénominateurs caractéristiques de Bose--Einstein et de Fermi--Dirac.
La condition bosonique sur le potentiel chimique garantit la
convergence des traces modales. Dans le régime thermodynamique, la
variation de l’entropie combinatoire sous contraintes d’énergie et de population
fournit une caractérisation complémentaire des mêmes lois.
Dans l’écran fini déclaré,
les dénombrements microcanoniques reproduisent par convolution les cardinalités
fermionique et bosonique au moyen de l’arithmétique entière. Les bornes du régime
dilué quantifient l’approximation commune de Maxwell--Boltzmann.

La loi de Planck résulte de la composition de trois données de la réalisation
électromagnétique libre : occupation bosonique à potentiel chimique nul,
densité de deux modes transversaux et énergie par excitation.
Le mode constant reste fixé comme fond. Les estimations
uniformes à l’origine et dans la queue spectrale justifient le passage des
sommes discrètes de la cavité aux densités continues. Les moments
de Bose déterminent les densités de nombre, d’énergie et d’entropie.
L’intégration angulaire du flux isotrope fournit la loi de
Stefan--Boltzmann ; l’extrémisation spectrale fournit les maxima
de Wien en fréquence et en longueur d’onde.

La constante de Stefan--Boltzmann est le coefficient de la composition
entre occupation bosonique, densité de modes et projection angulaire du
flux. Son expression en unités de calendrier relie l’action,
la durée et l’information tritique. Les constantes de Wien caractérisent
les points stationnaires de deux densités reliées par un changement de
mesure, dont le jacobien explique la différence entre leurs maxima.
Ces constructions distinguent les coefficients structurels, les
unités dimensionnelles et les grandeurs expérimentales auxquelles
ils sont comparés.

La constante de Boltzmann agit comme transducteur entre l’information
sans dimension et l’entropie dimensionnelle. L’énergie de décision provient
de l’action et de la période ; la section thermique spécifiée
individualise le transducteur et la température, avec leurs lois de changement
de base. La relation thermique du calendrier conserve l’unité
d’action, l’unité temporelle et la paire Boltzmann--température.
L’origine de l’énergie, la mesure des continuations et leur lecture
thermique sont ainsi reliées par des opérations de types
différenciés.

La chaîne renouvellement--quotient--mémoire--lacet réalise ce lien. Le renouvellement
$(20,24,25,23)$ et les quotients tritiques produisent le porteur gradué ;
la mémoire positive conserve $N-L$ ; le lacet $8{:}1$ impose
$\nu_W=11/9$ et $r=1/10$ ; et le regroupement $j=3n+b$ produit
$q=1/1000$ sans perdre le résidu. Avec les seuils $23,26,29$ et les
multiplicités $1,380,649$, on obtient exactement
$1{,}380649\times10^{-23}$. Les courants ascendants et descendants
imposent le déphasage sectoriel nul, tandis que le marqueur restitue la
partition après normalisation. L’équivalence centrée conserve la
réponse thermique et l’origine conserve le coût absolu en aire, en masse et en
rayon. Enfin,
$G_a k_B\Theta_{{\rm clk},a}\log3=c^4L_\alpha$ et l’identité du
transducteur composent la sortie spectrale avec la section dimensionnelle.

La dynamique qui précède l’occupation dispose d’une réalisation
unitaire explicite. Le Hamiltonien du calendrier a une exponentielle
qui reproduit exactement le décalage ; son action variationnelle donne
l’équation de Schrödinger correspondante. La formulation générale
maintient la continuité forte et le domaine du générateur autoadjoint.
Dans la réalisation périodique finie, le propagateur de Floquet reproduit
le décalage de \(108\) états. Pour les états de base
\(\rho_{j_0}=|j_0\rangle\langle j_0|\), le vecteur des espérances des
deux observables d’ordre a une période primitive de \(108\) cycles
de l’excitation périodique.
La conjugaison unitaire transporte
conjointement l’état, l’observable et le propagateur. Pour les perturbations
générales, l’estimation spectrale et la borne temporelle de
l’erreur sont conservées, en distinguant la périodicité, la covariance et la stabilité quantifiée.

\Needspace{11\baselineskip}
L’interprétation de l’entropie dépend de l’objet sur lequel elle est
évaluée. Pour l’état fermionique quasi libre, conservant le nombre
et invariant sous les transformations de jauge en dimension finie,
le corrélateur à un corps \(0<C<I\) et son générateur modulaire \(K_1\)
se déterminent réciproquement au moyen de la loi logistique.
Les entropies d’occupation, d’un opérateur densité et d’une
réduction intriquée sont reliées par des applications explicites et
conservent leurs domaines respectifs. L’existence de la limite
thermodynamique exige l’indépendance par rapport à la suite cofinale ;
la condensation est caractérisée par la saturation du secteur excité.

La mesure stationnaire de l’enveloppe surfacique--volumétrique est fixée
par l’incidence et la normalisation. Son taux maximal \(\log\varphi\)
s’accompagne d’une identité de divergence relative qui quantifie le
déficit de toute autre mesure stationnaire. Le rapport des poids
\(\varphi^{-2}\) appartient à cette pondération des histoires ; l’orbite
chronologique conserve sa propre répartition temporelle. La distinction permet
d’identifier la dynamique qui produit chaque distribution et la signification
de l’entropie correspondante.

L’aire totale et l’opérateur modulaire fournissent deux observables
complémentaires : la première conserve l’occupation conjointe et le second
pondère différemment les deux incidences lorsque \(\Delta_{\rm mod}\ne0\).
Sous cette condition et avec une normalisation connue, leur lecture simultanée
récupère les occupations sectorielles ; pour \(\Delta_{\rm mod}=0\),
le Hamiltonien modulaire est scalaire et cette récupération n’est pas possible.
L’identité finie exprime la variation d’entropie comme la somme
pondérée des variations surfaciques moins un déficit d’entropie relative
non négatif. L’inversion des occupations
conserve l’entropie et transforme l’aire en son complémentaire ; les deux
propriétés sont déterminées sans identifier l’égalité entropique
à l’égalité de structure.

La conservation de la mémoire soutient ces relations. Le noyau
arithmétique maintient les deux évaluations de somme et de produit, leurs
quotients et l’histoire des transitions, en distinguant l’état
transporté de ses observables. Le couplage bilatéral conserve
la norme et récupère l’entrée à partir des deux canaux ; l’invariance
d’autres observables obéit à la condition de commutation établie.
La récupération de l’état et la préservation d’une lecture particulière
sont donc des propriétés déterminées par leurs opérateurs respectifs.

La réalisation de Fourier du registre positionnel complète cette thèse de
conservation par deux bornes exactes. Tout état non nul sur
\(\Lambda=(\mathbb Z/27\mathbb Z)^2\) satisfait un produit de supports
au moins égal à \(729\), et tout état unitaire satisfait
\(H_{\rm pos}+H_F\ge\log729\). Les indicatrices normalisées de quatre
sous-groupes rectangulaires saturent simultanément les deux bornes. La preuve
entropique s’obtient par interpolation de Riesz--Thorin et différentiation
en \(p=2\). Ces entropies de distribution conservent leur domaine propre,
distinct des domaines de l’intrication, de von Neumann et de la thermodynamique.
Les lecteurs dodécaphasiques complémentaires prouvent, en outre, que deux
mots ayant un histogramme identique peuvent conserver des corrélations ordonnées
distinctes.

L’observateur premier des vacances retient une autre information que le dénombrement
marginal écarte : il combine \(z_p\) avec le poids \(\log p\) et produit
\(R_{\rm vac}(X)\). Sa résolution modale est une identité à convergence
symétrique ou de Fejér ; les paires conjuguées garantissent sa réalité. Les six
échelles conservées et les trente modes de \(X=10^8\) sont des témoins finis,
avec reproduction indépendante jusqu’à \(10^6\). La promotion asymptotique
requiert des bornes modales uniformes ; le crible conventionnel intervient ensuite
comme lecteur de vérification.

La densité photonique est en outre reliée, sans changer son générateur, au
niveau régularisé \(A_2^{\mathrm{BG}}\) : l’équation fonctionnelle donne
\(\mathcal Z_3(3)=4\pi_{\mathrm{HMT}}^2\log A_2^{\mathrm{BG}}\), et sa
substitution dans les moments thermiques produit exactement la densité
sans dimension de nombre photonique et les rapports moyens d’énergie et d’entropie.
Cette composition est la réciproque photonique du développement spectral
Bendersky--Glaisher ; elle conserve séparées la généalogie HMT de \(\pi\), la
régularisation spectrale et la réalisation électromagnétique.

La distribution complète de l’aire et l’entropie admettent des préparations ayant
des énergies distinctes. La section~\ref{ins29:v:orientacion} situe cette
différence dans l’information intrafibre : l’échange sectoriel conserve
l’aire, et son inverse atteint le maximum d’énergie récupérable pour le
Hamiltonien spécifié. Le restituteur conditionné caractérise en outre
le minimum d’énergie libre à marginale fixe.


Le résultat conjoint est une chaîne explicite de construction :
la trajectoire et sa parité déterminent les complétions d’occupation ;
l’équilibre détermine Bose--Einstein et Fermi--Dirac ; et l’occupation
bosonique, réalisée sur les modes électromagnétiques, détermine
Planck, Stefan--Boltzmann et Wien. L’action, le calendrier et la
transduction thermique relient leurs échelles, tandis que les applications de mémoire
et de réduction précisent l’information conservée. Les lois statistiques
et radiatives sont reliées par cette structure, avec les
conditions de chaque composition exprimées dans leurs domaines propres.

\par\endgroup
